\documentclass{tufte-handout}

\usepackage{style}

\begin{document}
\bibliography{biblio.bib} 
\tma{01}

%Transscribed from handwritten notes
\begin{question}

    \qpart

    \begin{center}
\textbf{MST125  TMA 01 Question 1}\\
Paul Allen Y362220X\\
\LaTeX  compiled with \TeX Live 2019
    \end{center}

\qpart
The distance between \( A\rb{x_1,y_1} \) and \( B\rb{x_2,y_2} \)
is given by 
\[ AB = \sqrt{\rb{x_2 - x_1}^2 + \rb{y_2 - y_1}^2} \]

Hence the distance between \( A\rb{1,4} \) and \( B\rb{3,-1} \) is
\begin{align*}
    AB &= \sqrt{\rb{1-3}^2 + \rb{4-(-1)}^2} \\
    &= \sqrt{\rb{-2}^2 + \rb{5}^2} \\
    &= \sqrt{4 + 25} \\
    &= 29
\end{align*}

\qpart
The gradient of the line thriough \( \rb{x_{1},y_{1}} \)\\
and \( x_{2},y_{2} \) is given by
\[ m = \frac{y_{2} -y_{1} }{x_{2} - x_{1}} \]

Hence the gradient of the line through\\
\( \rb{1,4} \) and \( \rb{3,-1} \) is
\[ m = \frac{\rb{-1} - 4}{3 - 1} = -\frac{5}{2} \]

\qpart
The gradient of the line is \( -\frac{5}{2} \) Hence\\
\( \tan\alpha = -\frac{5}{2} \).

Let \( \phi \) be the acute angle that the line\\
makes with the negative direction of the \( x-axis \).

\begin{align*}
\stext{Then}\\
    \tan\phi &= \frac{5}{2}, \\
\stext{so}\\
    \phi &= \taninv\rb{\frac{5}{2}} = 1.950\ldots \\
\stext{Hence}\\
    \alpha &= \pi - 1.190\ldots = 1.951\ldots\\
\stext{Therefore the angle \( \alpha \) is \( \SI{1.95}{\rad} \)}\\
\stext{(to 2 d.p.)}.
\end{align*}

\end{question}

%Question two

\begin{question}
%Least residue of 405^10 \pmod{16}
\qpart

\begin{align*}
    405 &\equiv 5 \pmod{16}\\
\stext{It follows that}
    405^{10} &\equiv 5^{10} \pmod{16}\\[8pt]
    5^{2} = 25 &\equiv 9 \pmod{16}\\[8pt]
    5^{4} = 25^{2} &\equiv 9^{2} = 81 \equiv 1 \pmod{16}\\
\stext{\( 5^{10} = 5^{8}\cdot{5^{2}}\)}
    5^{8} = \rb{5^4}^{2} &\equiv 1^{2} = 1 \pmod{16}\\
\stext{So}
    5^{10} &\equiv 1\cdot{5^{2}} \equiv 9 \pmod{16}\\
\stext{Hence}
    405^{10} &\equiv 9 \pmod{16}
\end{align*}

\vspace{5cm}

\marginnote{\[ a^{p-1} \equiv 1 \pmod{p}\]}
%Fremats little theorem to find least residue of 13^{494} \pmod{83}
\qpart

\begin{align*}
\stext{Using Fermat's Little Theorem, as \( 83 \) is prime}\\
    13^{82} &\equiv 1 \pmod{83}\\[8pt]
\stext{That implies}\\
    13^{164} &\equiv \rb{13^{82}}^2 \equiv 1^2 \equiv 1 \pmod{83}\\[8pt]
    13^{328} &\equiv \rb{13^{82}}^4 \equiv 1^4 \equiv 1 \pmod{83}\\[8pt]
    13^{492} &\equiv \rb{13^{82}}^6 \equiv 1^6 \equiv 1 \pmod{83}\\[8pt]
\stext{Hence, for \( 13^{494} \)}\\
    13^{494} &= 13^{492}\cdot{13^{2}} \equiv 1\cdot{13^{2}} \pmod{83}\\[8pt]
\stext{So}\\
    &= 169 \pmod{83}\\[8pt]
    &= 3 \pmod{83}\\
\stext{Hence the least residue of \( 13^{494} \pmod{83} \) is \( 3 \)}
\end{align*}

\vspace{3cm}

%Fermats little theorem to find 13^{6p-4} \equiv b \pmod{p}
\qpart

Let \( p \) be a prime \( p \neq 13 \) find a positive integer \( b \) such that 
\[ 13^{6p-4} \equiv b \pmod{5} \]

\begin{align*}
\stext{Let p=5}\\
    13^{\rb{6\rb{5} -4}} &\equiv b \pmod{5}\\[8pt]
    13^{26} &\equiv b \pmod{5}\\[8pt]
    13^{2} = 169 &\equiv 4 \pmod{5}\\
\stext{Using Fermat's Little Theorem}\\
    13^{4} &\equiv 1 \pmod{5}\\[8pt]
    13^{16} &equiv 1 \pmod{5}\\[8pt]
    13^{26} = 13^{16}\cdot{13^{4}}\cdot{13^{2}} &\equiv 1\cdot{1}\cdot{4} \equiv 4 \pmod{5}
\end{align*}

\end{question}

\clearpage
%Question three

\begin{question}
    
\qpart
% Find the multiplicative inverse of 29 \pmod{80} and hence solve 29x \equiv 41\pmod{80}
\qsubpart
\begin{align*}
\stext{Find the multiplicative inverse of \( 29 \pmod{80} \) using Euclid's algorithm}\\
    \circled{80} &= 2\cdot{\circled{29}} + \circled{22}\\[8pt]
    \circled{29} &= 1\cdot{\circled{22}} + \circled{7}\\[8pt]
    \circled{22} &= 3\cdot{\circled{7}} + \circled{1}\\[8pt]
    \circled{7} &= 1\cdot{\circled{7}} + 0\\
\stext{As the second to last remainder is \( 1 \) (Thus are co-prime) we can work backwards to find the multiplicative inverse.}\\
\stext{Backward substitution gives}\\
    \circled{1} &= 22 - 3\cdot{7}\\
    &= 22 - 3\cdot{\rb{\circled{29} - 22}}\\
    &= 22 - 3\cdot{\circled{29}} + 3\cdot{22}\\
    &= 4\cdot{22} - 3\cdot{\circled{29}}\\
    &= 4\rb{\circled{80} - 2\cdot{\circled{29}}} - 3\cdot{\circled{29}}\\
    &= 4\cdot{\circled{80}} - 8\cdot{\circled{29}} - 3\cdot{\circled{29}}\\
    &= 4\cdot{\circled{80}} - 11\cdot{\circled{29}}\\
\stext{Hence}
    -11 &\equiv 69 \pmod{80}\\
\stext{Hence the multiplicative inverse of \( 29 \pmod{80} \) is \( 69 \)}\\
\stext{To solve \( 29x \equiv 41 \pmod{80} \) multiply both sides by the multiplicative inverse}\\
    69\cdot{29x} &\equiv 69\cdot{41} \pmod{80}\\[8pt]
\stext{Since}\\
    69\cdot29 &\equiv 1 \pmod{80}\\
\stext{We have}\\
    x &\equiv 69\cdot{41} \pmod{80}\\
    x &\equiv 2829 \pmod{80}\\
    x &\equiv 29 \pmod{80}\\
\stext{Hence the solution is \( x \equiv 29 \pmod{80} \)}
\end{align*}

\clearpage

\qsubpart

\begin{align*}
\stext{Given}
    12x &\equiv 16 \pmod{54}
\stext{As \( 12 \) and \( 54 \) are not co-prime, we cannot use Euclid's algorithm to find the multiplicative inverse.}\\
\stext{Instead, we can divide both sides by the greatest common divisor of \( 12 \) and \( 54 \)}\\
\stext{The greatest common divisor of \( 12 \) and \( 54 \) is \( 6 \)}\\
\stext{But, \( 6 \) does not divide \( 16 \)}\\
\stext{Hence, the equation has no solution.}
\end{align*}

\vspace{3cm}

\qsubpart

\begin{align*}
\stext{Given}
    12x &= 18 \pmod{54}\\
\stext{As \( 12 \) and \( 54 \) are not co-prime, we cannot use Euclid's algorithm to find the multiplicative inverse.}\\
\stext{Instead, we can divide both sides by the greatest common divisor of \( 12 \) and \( 54 \)}\\
\stext{The greatest common divisor of \( 12 \) and \( 54 \) is \( 6 \)}\\
\stext{Dividing both sides by \( 6 \) gives}\\
    2x &\equiv 3 \pmod{9}\\
\stext{As the numbers are small we can try all possible values of \( x \) to find the solution.}\\
    2 \times 1 &\equiv 2 \pmod{9}\\
    2 \times 2 &\equiv 4 \pmod{9}\\
    2 \times 3 &\equiv 6 \pmod{9}\\
    2 \times 4 &\equiv 8 \pmod{9}\\
    2 \times 5 &\equiv 1 \pmod{9}\\
    2 \times 6 &\equiv 3 \pmod{9}\\
\stext{Hence the solution is \( x \equiv 6 \pmod{9} \)}
\end{align*}

\vspace{3cm}

\qpart

\[ E\rb{x} \equiv 19x - 5\pmod{26} \]

\begin{tabular}{|c|c|c|c|c|c|c|c|c|c|c|c|c|}
\hline
A& B & C & D & E & F & G & H & I & J & K & L & M\\
0& 1 & 2 & 3 & 4 & 5 & 6 & 7 & 8 & 9 & 10 & 11 & 12\\
\hline
N & O & P & Q & R & S & T & U & V & W & X & Y & Z\\
13& 14 & 15 & 16 & 17 & 18 & 19 & 20 & 21 & 22 & 23 & 24 & 25\\
\hline
\end{tabular}

So we can write the rule for deciphering as
\[ D\rb{y} \equiv v\rb{y + 5}\pmod{26} \]
\hfill (\text{Were \(v\) is the multiplicative inverse of } \( 19 \pmod{26} \))

\qsubpart
\begin{align*}
\stext{Using the Euclidean Algorithm to find the multiplicative inverse of \( 19 \pmod{26} \)}\\
    26 &= 1\cdot{19} + 7\\
    19 &= 2\cdot{7} + 5\\
    7 &= 1\cdot{5} + 2\\
    5 &= 2\cdot{2} + 1\\
    2 &= 2\cdot{1} + 0\\
\stext{As the second to last remainder is \( 1 \) (Thus co-prime)}\\
\stext{Backward substitution gives}\\
    1 &= 5 - 2\cdot{2}\\
    &= 5 - 2\cdot{\rb{7 - 5}}\\
    &= 5 - 2\cdot{7} + 2\cdot{5}\\
    &= 5 - 2\cdot{7} + 2\cdot{\rb{19 - 2\cdot{7}}}\\
    &= 5 - 2\cdot{7} + 2\cdot{19} - 4\cdot{7}\\
    &= 5 - 6\cdot{7} + 2\cdot{19}\\
    &= 5 - 6\cdot{26 - 19} + 2\cdot{19}\\
    &= 5 - 6\cdot{26} + 6\cdot{19} + 2\cdot{19}\\
    &= 5 - 6\cdot{26} + 8\cdot{19}\\
    &= 19 - 2\cdot{7} + 8\cdot{19}\\
    &= 19 - 2\cdot{\rb{26 - 19}} + 8\cdot{19}\\
    &= 19 - 2\cdot{26} + 2\cdot{19} + 8\cdot{19}\\
    &= 11\cdot{19} - 2\cdot{26}\\
\stext{Hence the multiplicative inverse of \( 19 \pmod{26} \) is \( 11 \)}\\
\snote{As required}
\end{align*}

\vspace{3cm}

\qsubpart

Hence we can use
\[ D\rb{y} \equiv 11\rb{y + 5}\pmod{26} \]
to decipher the message \( 14, 11, 5 \)

\begin{align*}
D\rb{14} &\equiv 11\rb{14 + 5} \pmod{26}\\
    &\equiv 11\cdot{19} \pmod{26}\\
    &\equiv 209 \pmod{26}\\
    &\equiv 1 \pmod{26}\\
    &\equiv B\\
D\rb{11} &\equiv 11\rb{11 + 5} \pmod{26}\\
    &\equiv 11\cdot{16} \pmod{26}\\
    &\equiv 176 \pmod{26}\\
    &\equiv 20 \pmod{26}\\
    &\equiv U\\
D\rb{5} &\equiv 11\rb{5 + 5} \pmod{26}\\
    &\equiv 11\cdot{10} \pmod{26}\\
    &\equiv 110 \pmod{26}\\
    &\equiv 6 \pmod{26}\\
    &\equiv G
\end{align*}

Hence the message is \( BUG \)

\end{question}

\clearpage

%Question four
\begin{question}

Given the elipse
\[ 4x^{2} + 25y^{2} - 9 = 0 \]

\qpart
\qsubpart

\begin{align*}
\stext{When \( x=0 \)}
    25y^{2} &= 9\\[8pt]      
    y^{2} &= \frac{9}{25}\\[8pt]
    y &= \pm\frac{3}{5}
\stext{When \( y=0 \)}
    4x^{2} &= 9\\[8pt]
    x^{2} &= \frac{9}{4}\\[8pt]
    x &= \pm\frac{3}{2}[8pt]
\stext{Hence the vertices are \[ \rb{\pm\frac{3}{2},0} \], \[ \rb{0,\pm\frac{3}{5}} \] }
\end{align*}

\clearpage

\qsubpart

\marginnote{ \[ e = \sqrt{1 - \frac{b^{2}}{a^{2}}} \]}

To figure out the eccentricity of the elipse, we can use,

\begin{align*}
\stext{Using the values of \( a \) and \( b \) from the equation of the elipse}\\
    e &= \sqrt{1 - \frac{\rb{\frac{3}{5}}^{2}}{\rb{\frac{3}{2}}^{2}}}\\[8pt]
    &= \sqrt{1 - \frac{9}{25}\cdot{\frac{4}{9}}}\\[8pt]
    &= \sqrt{1 - \frac{4}{25}}\\[8pt]
    &= \sqrt{\frac{21}{25}}\\[8pt]
    &= \frac{\sqrt{21}}{5}
\end{align*}

\marginnote{ \[ \text{The foci of the elipse are at } \rb{\pm ae,0} \]}

To find the foci of the elipse, we can use,

\begin{align*}
\stext{Using the values above, we have}\\
    foci &= \rb{\pm\frac{3}{2}\cdot{\frac{\sqrt{21}}{5}},0}\\[8pt]
    &= \rb{\pm\frac{3\sqrt{21}}{10},0}
\end{align*}

\marginnote{ The equation of the directrices is \[ x = \pm\frac{a}{e} \]}

\begin{align*}
\stext{Using the values of \( a \) and \( e \) from the equation of the elipse, the directories are}\\
    x &= \pm\frac{\frac{3}{2}}{\frac{\sqrt{21}}{5}}\\[8pt]
    &= \pm\frac{5\sqrt{21}}{14}\\[8pt]
\end{align*}

\clearpage

\qsubpart

\includegraphics[scale=0.7]{question_4.pdf}

\vspace{3cm}

\qsubpart

Using the standard form of the equation of an elipse

\begin{align*}
    4^{2} + 25^{2} - 9 &= 0\\
    4x^{2} + 25y^{2} &= 9\\
\stext{Dividing by \( 9 \)}\\
    \frac{4x^{2}}{9} + \frac{25y^{2}}{9} &= 1\\
\stext{Hence the equation of the elipse is}\\
    \frac{x^{2}}{\rb{\frac{3}{2}}^{2}} + \frac{y^{2}}{\rb{\frac{3}{5}}^{2}} &= 1
\end{align*}

Therefore the parameterisation of the elispe in the third quadrand is

\[ x = \frac{3}{2}\cos\rb{t}, y = \frac{3}{5}\sin\rb{t} \]
\hfill Where \( \pi < t < \frac{3}{2}\pi \)


\qsubpart

When this part of the conic is moved 4 units to the right and 2 units down, it would
have parameterisation of
\[ x = 4 + \frac{2}{3}\cos\rb{t}, y = -2 + \frac{5}{3}\sin\rb{t} \]

\vspace{3cm}

\qpart
\qsubpart

\marginnote{General formula for a conic \[ Ax^{2} + Bxy +Cy^{2} + Dx + Ey + F = 0 \]}

Using the general formula for the equation for a conic, and my OU student number being
Y362220x, hense the last nono zero digit is \( 2 \) therefore \( B = 2 \)

\begin{align*}
    -2x^{2} + Bxy + 3y^{2} - x + 3y - 5 &= 0\\
    -2x^{2} + 2xy + 3y^{2} - x + 3y - 5 &= 0
\end{align*}

Therefore \( A=-2, B=2 \text{ and } c=3 \)

\begin{align*}
\stext{evaluating \( B^{2} - 4AC \)}
    2^{2} - 4\cdot{-2}\cdot{3} &= 4 + 24\\
    &= 28\\
    28 > 0\\
\stext{Hence the equation represents a hyperbola}
\end{align*}

\clearpage

\qsubpart

\includepdf[pages=-]{question_4_b.pdf}
\includepdf[pages=-]{question_4_b_ii.pdf}

\end{question}

\clearpage

%Question five
\begin{question}

\qpart

The position of \textbf{A} is \( \rb{t-3,3t+8} \)

\begin{align*}
    x &= t-3\\
    t &= x+3\\
    \stext{Hence the equation fo the line is given by}\\
    y &= 3\rb{x+3} + 8\\
    &= 3x + 9 + 8\\
    &= 3x + 17
\end{align*}

\vspace{3cm}

\qpart

\marginnote{The distance between two points is given by \[ d^2 = \rb{x_{2}-x_{1}}^{2} + \rb{y_{2}-y_{1}}^{2} \]}

The position of \textbf{B} is \( \rb{2t-5,4t+5} \)

\begin{align*}
\stext{Substituting our position of \textbf{A} and \textbf{B} into the distance formula}\\
    d^{2} &= \rb{2t-5 - \rb{t - 3 }}^{2} + \rb{4t + 5 - \rb{3t + 8}}^{2}\\[8pt]
    &= \rb{t - 2}^{2} + \rb{t - 3}^{2}\\[8pt]
    &= t^{2} - 4t + 4 + t^{2} - 6t + 9\\[8pt]
    &= 2t^{2} - 10t + 13
\snote{as required}
\end{align*}

\clearpage

\qpart

\marginnote[2cm]{The first mistake that was made was when completing the square, the expression should
have been \[ = 2\rb{t - \frac{5}{2} - \frac{25}{4}} +13 \] \[ = 2\rb{t - \frac{5}{2}} -\frac{25}{2} +13  \].
}

\marginnote[1cm]{If we contiue from this error, the second mistake was when solving for the distance, the expression should have been
\[ d^{2} = \frac{27}{4} \] hence \[ d=\sqrt{\frac{27}{4}} = 2.59\ldots \]}
   
\begin{align*}
        d^2 &= 2t^{2} - 10t + 13\\[8pt]
        &= 2\rb{t^{2} -5t} +13\\[8pt]
        &= 2\rb{t - \frac{5}{2}}^2 - \frac{25}{4} +13\\[8pt]
\stext{Note 1}
        &= 2\rb{t - \frac{5}{2}}^2 + \frac{27}{4}\\[8pt]
    \stext{The mimimum value of \( d^{2} \) occurs when \( t=\frac{5}{2}=2.5\). 
    Hence the minium distance is \( \frac{27}{4} \)}
\end{align*}

\vspace{5cm}

My soloution is

\begin{align*}
    d^{2} &= 2t^{2} - 10t + 13\\
    &= 2\rb{t^{2} - 5t} +13\\
\stext{Completing the square}\\
    &= 2\rb{\rb{t - \frac{5}{2}}^{2} - \frac{25}{4}} + 13\\[8pt]
    \stext{simplyfying}\\
    &= 2\rb{t - \frac{5}{2}}^{2} - \frac{25}{2} + 13\\[8pt]
    &= 2\rb{t - \frac{5}{2}}^{2} + \frac{1}{2}\\[8pt]
\stext{The minimum value of \( d^{2} \) occurs when \( t = \frac{5}{2} \)}\\
    &= \frac{1}{2}\\[8pt]
    &= 0.5\\
\snote{to 2 s.f}
\stext{Hence}\\
    d &= \sqrt{\frac{1}{2}}\\[8pt]
    &= \frac{\sqrt{2}}{2}\\[8pt]
    &= 0.7071\ldots\\
    &= 0.71\\
\snote{to 2 s.f}
\end{align*}


\end{question}

\clearpage

%Question six
\begin{question}

\qpart
\qsubpart

\begin{minipage}{0.7\textwidth}
    \begin{tikzpicture}
        \tkzDefPoint(0,0){P}
        \tkzDefPoint(0,-2){W}
        \tkzDefPoint(0,2){T}
        \tkzDrawSegment[->](P,W)
        \tkzDrawSegment[->](P,T)
        \tkzDrawPoint(P)
        \tkzLabelPoint[below](W){\( W_{potatoes} \)}
        \tkzLabelPoint[above](T){\( T_{pulley} \)}
        \tkzLabelPoint[right](P){\( P \)}
    \end{tikzpicture}
\end{minipage}
\begin{minipage}{0.3\textwidth}
\begin{tikzpicture}
    \tkzDefPoint(0,0){o}
    \tkzDefShiftPoint[o](0:1){i}
    \tkzDefShiftPoint[o](90:1){j}
    \tkzDrawSegments[-{stealth}](o,i o,j)
    \tkzLabelPoint[below](i){\textbf{i'}}
    \tkzLabelPoint[left](j){\textbf{'j}}
\end{tikzpicture}
\end{minipage}

As the system is in equalibrium, 
\begin{align*}
    T + W &= 0\\
\stext{Or}\\
T &= -W\\
\stext{With}\\
    W &= 10g\\
\snote{where \( g \) is the acceleration due to gravity}\\
    T &= -10g\\
\end{align*}

\qsubpart

\begin{align*}
\stext{In component from}\\
    T &= 10g\textbf{j'}\\
\stext{and}\\
    W &= -10g\textbf{j'}\\
\end{align*}

\qpart

\qsubpart

Force diagram of a mass on a slope.

\begin{minipage}{0.7\textwidth}
\begin{tikzpicture}
    \tkzDefPoint(0,0){M}
    \tkzDefShiftPoint[M](270:2){W}
    \tkzDefShiftPoint[M](45:2){N}
    \tkzDefShiftPoint[M](105:2){T}
    \tkzDefShiftPoint[M](315:2){F}
    \tkzDefShiftPoint[M](135:2){S}
    \tkzDrawSegments[->](M,W M,N M,T M,F) 
    \tkzDrawLine[dashed](M,S)
    \tkzDrawPoint(M)
    \tkzLabelPoint[below](W){$W_{books}$}
    \tkzLabelPoint[right](M){$M$}
    \tkzLabelPoint[above](N){$N$}
    \tkzLabelPoint[above](T){$T_{potatoes}$}
    \tkzLabelPoint[below](F){$F$}
    \tkzMarkAngle(W,M,F)
    \tkzMarkAngle(T,M,S)
    \tkzLabelAngle[pos=0.8](W,M,F){\( \ang{45} \)}
    \tkzLabelAngle[pos=0.8](T,M,S){\( \ang{30} \)}
\end{tikzpicture}
\end{minipage}
\begin{minipage}{0.3\textwidth}
\begin{tikzpicture}
\tkzDefPoint(0,0){o}
\tkzDefShiftPoint[o](315:1){i}
\tkzDefShiftPoint[o](45:1){j}
\tkzDrawSegments[-{stealth}](o,i o,j) 
\tkzLabelPoints[below](i){\textbf{i}}
\tkzLabelPoints[left](j){\textbf{j}}
\end{tikzpicture}
\end{minipage}

\qsubpart

The forces acting on the box of books in component form are

\begin{align*}
    \mathbf{W_{books}} &= Mg\cos{45}\textbf{i} - Mg\sin{45}\textbf{j}\\[8pt]
    \mathbf{N} &= N\textbf{j}\\[8pt]
    \mathbf{T_{potatoes}} &= -10g\cos{30}\textbf{i} + 10g\sin{30}\textbf{j}\\[8pt]
    \mathbf{F} &= F\textbf{i}
\snote{Were \( g \) is the acceleration due to gravity}
\end{align*}

\qsubpart

As the whole system is in Equlilibrium

\[ \mathbf{T_{potatoes}} + \mathbf{W_{books}} + N + F = 0 \]

\qsubpart

\begin{align*}
\stext{Resolving this is the \( \mathbf{i}-direction \)}
    -10g\cos{30} + Mg\cos{45} + F &= 0\\
    F &= 10g\cos{30} - Mg\cos{45}\\
\stext{When \( \SI{12}{\kilo\gram} \)}
    F &= 10\cdot{9.8}\cos{30} - 12\cdot{9.8}\cos{45}\\
    &= 49\sqrt{3} - 83.155\ldots\\
    &= -1.7\si{\newton}\\
\snote{to 2 s.f}\\
\stext{when \( \SI{13}{\kilo\gram} \)}
    F &= 10\cdot{9.8}\cos{30} - 13\cdot{9.8}\cos{45}\\
    &= 49\sqrt{3} - 90.085\ldots\\
    &= 5.2\si{\newton}\\
\snote{to 2 s.f}
\end{align*}

So when the mass of the books is \( \SI{12}{\kilo\gram} \) the force imparted by friction is 
\( -1.7\si{\newton} \) and is running up the slope in the \( \mathbf{i}-direction \).

When the mass of the books is \( \SI{13}{\kilo\gram} \) the force imparted by friction is
\( 5.2\si{\newton} \) and is running down the slope in the \( -\mathbf{i}-direction \).

\end{question}

%Question 7

\begin{question}

\qpart

\includegraphics[scale=0.6]{question_7.jpg}

\qpart

Based on my work in Unit 5 and the practice quiz, I feel confident in resolving forces in a system 
and finding the resultant force acting on an object. My prior experience with mechanics has definitely 
helped. However, I still struggle with keeping track of signs when resolving forces. I often miss 
a negative sign, especially when transcribing my handwritten work into \LaTeX{}. This is partly due to my 
dyslexia, which makes spotting transcription errors difficult. While I can catch these mistakes when 
I print my work, it’s something I’ll need to be especially careful of during exams.

\vspace{3cm}

\qpart

To improve my understanding of Unit 5, I plan to revisit my previous Level 2 mechanics material, 
which I believe will reinforce the current content. I also aim to practise more problem-solving 
to strengthen my intuition and accuracy, particularly in sign conventions. Doing this will help 
me reduce mistakes and improve confidence. Additionally, I’ll review my printed drafts more 
carefully, as this method helps me catch errors more effectively. I scored well on the practice 
quiz overall, but these small mistakes could cost me marks under timed conditions. That’s why 
I think continued practice and review are essential.

\end{question}

\end{document}